\documentclass[11pt]{article}
\usepackage[margin=1in]{geometry}
\usepackage{amsmath,hyperref}
\begin{document}
\begin{center}\large A prescribed OH perturbation in Laughner et al. (2021)\end{center}
\noindent\textbf{Reference.} Laughner et al., \emph{PNAS} 118, e2109481118,
\href{https://doi.org/10.1073/pnas.2109481118}{doi:10.1073/pnas.2109481118}.
The main text, 19-page SI Appendix, team list, and both supplementary
workbooks (SI Tables~S2 and S4) were read.

\paragraph{Method.} This synthesis combines pandemic activity indicators,
observations, and chemistry experiments. Its methane analysis uses TCCON
columns at Park Falls and Lauder. Flag-zero data are averaged over 15 days,
deseasonalized, and compared with a robust linear 2016--2019 trend.
Growth rates use differences between three-month means multiplied by four.
Figure~6b averages monthly fractional deviations from each site's trend;
it is not a global mass-balance growth anomaly. A two-hemisphere methane--OH
box model includes emissions, oxidation, and interhemispheric transport.
Methane emissions follow a stabilized scenario. The plotted perturbation
compares runs with the same emissions, one having OH reduced by 3\% in 2020.

\paragraph{Findings and uncertainty.} Methane columns rise about 0.3\% above
their extrapolated trends in 2020, similar to the OH scenario. The authors
interpret this as support for a contribution from COVID-related NO$_x$
reductions and lower OH, while allowing emission changes. The 2--4\% OH
range comes from temporal variation in a chemistry-assimilation experiment:
2\% in March and 4\% in May. It is a proxy for uncertainty, not propagated
posterior uncertainty in an annual global OH estimate. The two 2021 lines
are extrapolations: one holds the January 2021 concentration difference
fixed, and one extends its final-six-month 2020 slope. They were not
calculated with the box model.

\paragraph{Interpretation for a 2019 baseline.} A uniform fractional
OH reduction $r$ gives the fixed-burden forcing
\begin{equation}
F_{{\rm OH},2020}=r\,m_{2019}/\tau_{{\rm OH},2019},
\end{equation}
in which $m$ is methane burden and $\tau_{\rm OH}$ is the OH-only lifetime.
We evaluate this conditional scenario using a NOAA burden proxy and an
11.1-yr benchmark, with 9--12-yr sensitivity. It is not a native flux
reconstruction or reproduction of the two-site observational anomaly.
The imposed 2020 forcing is zero again in 2021--2022; the accumulated
concentration response need not be zero.

\paragraph{Numerical data.} Dataset~S1 gives CO$_2$, CH$_4$, and NO$_x$
inventory series. Its CH$_4$ column ends in 2017; 2018--2020 are missing.
The equivalent-emission-year illustration uses extrapolation and an assumed
reduction, not an inferred 2020 global methane flux. Dataset~S2 gives monthly
ozone-production efficiencies for 20 cities, not methane sink fluxes. Both
workbooks are retained as numerical transcriptions. No 2019--2022 methane
budget is filled from their missing entries.
\end{document}
